% ============================================================
\chapter{Mesures de Radon et Th\'eor\`eme de Riesz}
\label{ch:radon_riesz}
% ============================================================

\section{Introduction}

Ce chapitre se situe \`a l'interface de la th\'eorie de la mesure et de la topologie.
Nous \'etudions les \emph{mesures de Radon}\index{mesure de Radon} sur les espaces
localement compacts et \'etablissons le c\'el\`ebre th\'eor\`eme de repr\'esentation de
Riesz\index{th\'eor\`eme de Riesz}, qui identifie les formes lin\'eaires positives sur
l'espace des fonctions continues \`a support compact avec les mesures de Radon.

\begin{intuition}[title=\textbf{Lien entre topologie et mesure}]
Les mesures de Radon sont des mesures qui \og respectent \fg la topologie~: elles
sont d\'etermin\'ees par leurs valeurs sur les compacts et les ouverts. Le th\'eor\`eme
de Riesz dit que toute mani\`ere \og raisonnable \fg d'int\'egrer des fonctions continues
\`a support compact provient d'une telle mesure. C'est le pont fondamental entre
analyse fonctionnelle et th\'eorie de la mesure.
\end{intuition}

% ------------------------------------------------------------
\section{Rappels sur les espaces localement compacts}
% ------------------------------------------------------------

\begin{definition}[Espace localement compact]
Un espace topologique s\'epar\'e $X$ est dit \emph{localement compact}\index{localement compact}
si tout point admet un voisinage compact, c'est-\`a-dire~: pour tout $x \in X$, il existe
un ouvert $U$ et un compact $K$ avec $x \in U \subset K$.
\end{definition}

\begin{exemple}
\begin{enumerate}
    \item $\R^n$ est localement compact.
    \item Tout espace compact est localement compact.
    \item Tout ouvert d'un espace localement compact est localement compact.
    \item Un espace de Banach de dimension infinie n'est pas localement compact.
\end{enumerate}
\end{exemple}

\begin{definition}[Support et fonctions \`a support compact]
Pour une fonction continue $f : X \to \R$, le \emph{support} de $f$ est
$\mathrm{supp}(f) = \overline{\{x \in X : f(x) \neq 0\}}$. On note
$C_c(X)$\index{$C_c(X)$} l'espace des fonctions continues \`a support compact.
\end{definition}

\begin{lemme}[Urysohn pour les espaces localement compacts]
\label{lem:urysohn}
Soit $X$ localement compact s\'epar\'e, $K$ compact et $U$ ouvert avec
$K \subset U$. Alors il existe $f \in C_c(X)$ avec $0 \leq f \leq 1$,
$f|_K = 1$ et $\mathrm{supp}(f) \subset U$.
\end{lemme}

\begin{proposition}[Partition de l'unit\'e]
\label{prop:partition_unite}
Soit $X$ localement compact s\'epar\'e et soit $K$ compact couvert par des ouverts
$U_1, \ldots, U_n$. Il existe des fonctions $\varphi_1, \ldots, \varphi_n \in C_c(X)$
avec $0 \leq \varphi_i \leq 1$, $\mathrm{supp}(\varphi_i) \subset U_i$, et
$\sum_{i=1}^n \varphi_i = 1$ sur $K$.
\end{proposition}

% ------------------------------------------------------------
\section{Mesures de Radon}
% ------------------------------------------------------------

\begin{definition}[Mesure bor\'elienne r\'eguli\`ere]
\label{def:reguliere}
Soit $X$ un espace localement compact s\'epar\'e. Une mesure positive $\mu$ sur
$\mathcal{B}(X)$ est dite~:
\begin{itemize}
    \item \emph{ext\'erieurement r\'eguli\`ere}\index{r\'egularit\'e ext\'erieure}
    si pour tout $B \in \mathcal{B}(X)$~:
    $\mu(B) = \inf\{\mu(U) : U \text{ ouvert},\, B \subset U\}$\,;
    \item \emph{int\'erieurement r\'eguli\`ere}\index{r\'egularit\'e int\'erieure}
    si pour tout ouvert $U$~:
    $\mu(U) = \sup\{\mu(K) : K \text{ compact},\, K \subset U\}$\,;
    \item \emph{r\'eguli\`ere} si elle est \`a la fois ext\'erieurement et
    int\'erieurement r\'eguli\`ere.
\end{itemize}
\end{definition}

\begin{definition}[Mesure de Radon]
\label{def:radon}
\index{mesure de Radon}
Une \emph{mesure de Radon} sur un espace localement compact s\'epar\'e $X$ est une
mesure bor\'elienne $\mu$ sur $X$ qui est~:
\begin{enumerate}
    \item[(i)] \emph{localement finie}~: pour tout $x \in X$, il existe un voisinage
    ouvert $U$ de $x$ avec $\mu(U) < \infty$\,;
    \item[(ii)] \emph{int\'erieurement r\'eguli\`ere sur les ouverts}~: pour tout
    ouvert $U$, $\mu(U) = \sup\{\mu(K) : K \text{ compact},\, K \subset U\}$\,;
    \item[(iii)] \emph{ext\'erieurement r\'eguli\`ere}~: pour tout bor\'elien $B$,
    $\mu(B) = \inf\{\mu(U) : U \text{ ouvert},\, B \subset U\}$.
\end{enumerate}
\end{definition}

\begin{exemple}
\begin{enumerate}
    \item La mesure de Lebesgue sur $\R^n$ est une mesure de Radon.
    \item Toute mesure de Dirac $\delta_a$ est une mesure de Radon.
    \item La mesure de comptage sur un espace discret est une mesure de Radon.
    \item La mesure de comptage sur $\R$ n'est \emph{pas} une mesure de Radon
    (elle n'est pas localement finie).
\end{enumerate}
\end{exemple}

\begin{proposition}
Toute mesure de Radon finie sur un espace localement compact s\'epar\'e est
r\'eguli\`ere (au sens int\'erieur sur \emph{tous} les bor\'eliens, pas seulement les
ouverts).
\end{proposition}

% ------------------------------------------------------------
\section{Fonctionnelles lin\'eaires positives}
% ------------------------------------------------------------

\begin{definition}[Fonctionnelle lin\'eaire positive]
Une forme lin\'eaire $\Lambda : C_c(X) \to \R$ est dite \emph{positive} si~:
\[
    f \geq 0 \implies \Lambda(f) \geq 0.
\]
\end{definition}

\begin{remarque}
Toute fonctionnelle lin\'eaire positive sur $C_c(X)$ est automatiquement continue
pour la topologie de la convergence uniforme sur les compacts (par la positivit\'e,
$\abs{\Lambda(f)} \leq \Lambda(\norm{f}_\infty \cdot \ind_K)$ o\`u $K$ contient le
support de $f$).
\end{remarque}

% ------------------------------------------------------------
\section{Th\'eor\`eme de repr\'esentation de Riesz}
% ------------------------------------------------------------

\begin{theoreme}[Riesz-Markov-Kakutani]
\label{thm:riesz}
\index{th\'eor\`eme de Riesz-Markov-Kakutani}
Soit $X$ un espace localement compact s\'epar\'e et soit $\Lambda : C_c(X) \to \R$
une fonctionnelle lin\'eaire positive. Alors il existe une unique mesure de Radon $\mu$
sur $X$ telle que~:
\[
    \Lambda(f) = \int_X f\, d\mu, \qquad \forall f \in C_c(X).
\]
\end{theoreme}

\begin{proof}[Id\'ee de la preuve]
La preuve se fait en plusieurs \'etapes~:

\emph{\'Etape 1.} D\'efinition de $\mu$ sur les ouverts. Pour $U$ ouvert, on pose~:
\[
    \mu(U) = \sup\{\Lambda(f) : f \in C_c(X),\, 0 \leq f \leq 1,\,
    \mathrm{supp}(f) \subset U\}.
\]

\emph{\'Etape 2.} Extension aux bor\'eliens par r\'egularit\'e ext\'erieure~:
\[
    \mu(B) = \inf\{\mu(U) : U \text{ ouvert},\, B \subset U\}.
\]

\emph{\'Etape 3.} V\'erification que $\mu$ est une mesure. On utilise le lemme
d'Urysohn et les partitions de l'unit\'e pour v\'erifier la $\sigma$-additivit\'e.

\emph{\'Etape 4.} V\'erification que $\Lambda(f) = \int f\, d\mu$. On approche $f$
par des fonctions \'etag\'ees et on utilise la r\'egularit\'e de $\mu$.

\emph{\'Etape 5.} Unicit\'e. Si $\mu'$ est une autre mesure de Radon avec
$\Lambda(f) = \int f\, d\mu'$, alors $\mu$ et $\mu'$ co\"incident sur les ouverts
(par le lemme d'Urysohn), donc sur les bor\'eliens (par r\'egularit\'e).
\end{proof}

\begin{attention}[title=\textbf{Conditions du th\'eor\`eme}]
L'espace $X$ doit \^etre localement compact et s\'epar\'e. La fonctionnelle doit \^etre
d\'efinie sur $C_c(X)$ (fonctions \`a support compact), pas sur $C_b(X)$ (fonctions
born\'ees). Pour des espaces plus g\'en\'eraux, on utilise d'autres versions du th\'eor\`eme.
\end{attention}

% ------------------------------------------------------------
\section{Support d'une mesure}
% ------------------------------------------------------------

\begin{definition}[Support d'une mesure]
\label{def:support_mesure}
\index{support d'une mesure}
Soit $\mu$ une mesure de Radon sur $X$ localement compact s\'epar\'e. Le
\emph{support} de $\mu$ est le plus petit ferm\'e $F$ tel que $\mu(X \setminus F) = 0$~:
\[
    \mathrm{supp}(\mu) = X \setminus \bigcup\{U \text{ ouvert} : \mu(U) = 0\}.
\]
\end{definition}

\begin{proposition}
Soit $\mu$ une mesure de Radon sur $X$. Alors~:
\begin{enumerate}
    \item $x \in \mathrm{supp}(\mu)$ si et seulement si $\mu(U) > 0$ pour tout
    voisinage ouvert $U$ de $x$.
    \item $\mathrm{supp}(\mu) = \varnothing$ si et seulement si $\mu = 0$.
    \item Si $f \in C_c(X)$ et $\mathrm{supp}(f) \cap \mathrm{supp}(\mu) = \varnothing$,
    alors $\int f\, d\mu = 0$.
\end{enumerate}
\end{proposition}

\begin{proof}
(1) $x \notin \mathrm{supp}(\mu)$ signifie que $x$ appartient \`a un ouvert $U$ de
mesure nulle, c'est-\`a-dire qu'il existe un voisinage ouvert de $x$ de mesure nulle.

(2) $\mathrm{supp}(\mu) = \varnothing$ signifie $\mu(X) = 0$ par r\'egularit\'e
int\'erieure (tout compact $K$ est contenu dans $X \setminus \mathrm{supp}(\mu)$,
qui est une r\'eunion d'ouverts de mesure nulle).

(3) $\mathrm{supp}(f)$ est compact et contenu dans l'ouvert
$X \setminus \mathrm{supp}(\mu)$ qui est de mesure nulle.
\end{proof}

\begin{exemple}
\begin{enumerate}
    \item $\mathrm{supp}(\delta_a) = \{a\}$.
    \item $\mathrm{supp}(\lambda|_{[0,1]}) = [0,1]$ o\`u $\lambda$ est la mesure
    de Lebesgue.
    \item La mesure de Cantor a pour support l'ensemble de Cantor.
\end{enumerate}
\end{exemple}

% ------------------------------------------------------------
\section{Th\'eor\`eme de Riesz pour les mesures sign\'ees}
% ------------------------------------------------------------

\begin{theoreme}
\label{thm:riesz_signe}
Soit $X$ localement compact s\'epar\'e. Toute forme lin\'eaire continue
$\Lambda : C_0(X) \to \R$ (o\`u $C_0(X)$ est l'espace des fonctions continues
s'annulant \`a l'infini, muni de la norme uniforme) s'\'ecrit de fa\c{c}on unique~:
\[
    \Lambda(f) = \int_X f\, d\mu
\]
o\`u $\mu$ est une mesure de Radon \emph{sign\'ee} r\'eguli\`ere avec
$\norm{\Lambda} = \abs{\mu}(X)$.
\end{theoreme}

\begin{formules}[title=\textbf{Correspondance de Riesz}]
\[
\boxed{
    \bigl(C_c(X)\bigr)^*_+ \;\longleftrightarrow\;
    \{\text{mesures de Radon positives sur } X\}
}
\]
\[
\boxed{
    \bigl(C_0(X)\bigr)^* \;\longleftrightarrow\;
    \{\text{mesures de Radon sign\'ees r\'eguli\`eres sur } X\}
}
\]
\end{formules}

% ------------------------------------------------------------
\section{Exercices}
% ------------------------------------------------------------

\begin{exercice}
Montrer que la mesure de Lebesgue sur $\R$ est une mesure de Radon.
\emph{Indication}~: v\'erifier les trois conditions de la d\'efinition~\ref{def:radon}.
\end{exercice}

\begin{exercice}
Soit $X = \R$ et $\Lambda(f) = f(0)$ pour $f \in C_c(\R)$. Identifier la mesure
de Radon $\mu$ telle que $\Lambda(f) = \int f\, d\mu$.
\end{exercice}

\begin{exercice}
Soit $\mu$ une mesure de Radon sur $\R^n$ et $K$ un compact avec $\mu(K) = 0$.
Montrer que pour tout $\varepsilon > 0$, il existe un ouvert $U \supset K$
avec $\mu(U) < \varepsilon$.
\end{exercice}

\begin{exercice}
Soit $X$ localement compact s\'epar\'e et $\mu$ une mesure de Radon finie. Montrer
que pour tout $B \in \mathcal{B}(X)$ et tout $\varepsilon > 0$, il existe un compact
$K \subset B$ et un ouvert $U \supset B$ avec $\mu(U \setminus K) < \varepsilon$.
\end{exercice}

\begin{exercice}[Support et densit\'e]
Soit $f : \R \to [0, +\infty)$ continue et $\mu(A) = \int_A f\, d\lambda$.
Montrer que $\mathrm{supp}(\mu) = \overline{\{x \in \R : f(x) > 0\}}$.
\end{exercice}

\begin{exercice}
Soit $X$ un espace localement compact s\'epar\'e et soient $\mu, \nu$ deux mesures
de Radon telles que $\int f\, d\mu = \int f\, d\nu$ pour toute $f \in C_c(X)$.
Montrer que $\mu = \nu$. Cet \'enonc\'e est-il vrai si l'on remplace $C_c(X)$ par
un sous-espace strict~?
\end{exercice}
